\documentclass{article}

%%% Package List

\usepackage{datetime} % Dates.
\usepackage{fancyhdr} % Headers.
\usepackage{lastpage} % Page References.

\usepackage{graphicx} % Inserting Images.
\usepackage{float} % Postitioning Images.

\usepackage{dirtytalk} % Speech & Quote Marks.
\usepackage{hyperref} % Hyperlinks and References.
\usepackage{xcolor} % Coloured Text.
\usepackage{enumitem} % Letter Bullets. 

\usepackage{amssymb} % Maths Symbols.
\usepackage{amsmath} % Maths Tools.

\usepackage{algpseudocode} % Pseudocode.
\usepackage{algorithm} % Algorithms.
\usepackage{minted} % Code Snippets.

\usepackage[margin=1in]{geometry} % Reduces the Page Margins.

%%%



%%% Page Variables
\newcommand{\thetitle}{Algorithms and Data Structures} % Insert the document title here.
\newcommand{\thesubtitle}{COMP1081: Durham University} % Insert the document subtitle here.
\newcommand{\theauthor}{Matthew Emerson} % Insert the name(s) of the author(s) here.
\newcommand{\thedate}{\monthyeardate\today} % Replace this if you want a static date.

\newcommand{\templateversion}{1.4.0} % Do not edit this variable.
%%%



%%% Custom Formats and Commands

\hypersetup{colorlinks=true, urlcolor=black, linkcolor=black}


% Date Formats.
\newdateformat{monthyeardate}
{%
    \monthname[\THEMONTH] \THEYEAR
}

\newdateformat{yeardate}
{%
    \THEYEAR
}


% Colouring and Spacing.
\newcommand{\colour}[2] {\color{#1}#2\color{black}}
\newcommand{\separate}[1] {\vspace{#1 pt}\noindent}


% Maths Number Sets.
\newcommand{\R}{\mathbb{R}}
\newcommand{\N}{\mathbb{N}}
\newcommand{\C}{\mathbb{C}}
\newcommand{\Q}{\mathbb{Q}}
\newcommand{\Z}{\mathbb{Z}}
\newcommand{\I}{\mathbb{I}}


% Custom Formats.
\newcommand{\definition}[2]{\textbf{Definition #1}: #2\label{def:#1}}
\newcommand{\example}[2]{\textbf{Example #1}: #2\label{ex:#1}}
\newcommand{\conjecture}[2]{\textbf{Conjecture #1}: #2\label{cnj:#1}}
\newcommand{\numberedentry}[3]{\textbf{#2 #1}: #3\label{entry:#1}}
\newcommand{\proof}[2]{\textit{Proof}: #2 \begin{flushright}\(\square\)\end{flushright}\label{prf:#1}}
\newcommand{\theorem}[3]{\textbf{Theorem #1}: #2 \par\separate{3}\proof{thm:#1}{#3}\label{thm:#1}}
\newcommand{\lemma}[3]{\textbf{Lemma #1}: #2 \par\separate{3}\proof{lem:#1}{#3}\label{lem:#1}}


% Advanced Formats.
\newcommand{\image}[3]
{\begin{figure}[H]
    \centering
    \includegraphics[width=#3\linewidth]{#1}
    \caption{#2}
    \label{fig:#2}
\end{figure}}

\newcommand{\algorithmBlock}[4]
{\begin{algorithm}[H]
    \caption{#1}
    \begin{algorithmic}
        \State \textbf{Input}: #2
        \State \textbf{Output}: #3
        \separate{5}#4
    \end{algorithmic}
    \label{alg:#1}
\end{algorithm}}

%%%



%%% Document Setup

% Title Page.
\title{\thetitle\\[0.5em]\large{\thesubtitle}}
\author{\theauthor}
\date{\thedate} 


% Custom Page Header.
\fancyhead{}
\fancyhead[R]{\thetitle\text{ - }\thesubtitle}

% Custom Page Footer.
\fancyfoot{}
\fancyfoot[L]{Written by \theauthor \\ 
\textit{\LaTeX \,Template v.\templateversion: Copyright \date{\yeardate\today}, Matthew Emerson}} % Do not edit.
\fancyfoot[R]{Page \thepage \, of\, \pageref{LastPage}}


% Setting Up Pages.
\pagestyle{empty}\begin{document}
\maketitle\thispagestyle{empty} % Title Page.
\newpage\tableofcontents % Contents Page.

\clearpage\setcounter{page}{1}\newpage % Excludes the title and contents page from page count.
\pagestyle{fancy}

%%%



%%% Document Content

\part{Fundamental Data Structures}\label{part:1.fds}

\section{Introduction}

\subsection{Basic Definitions}

\definition{1.1.1}{An \underline{algorithm} is a method or process that is followed to solve a problem. It must be correct, must be composed of a finite number of unambiguous steps, and must terminate.}

\separate{5}\definition{1.1.2}{A \underline{data structure} is a particular way of storing and organising data in a computer so that it can be used efficiently.}

\separate{5}\definition{1.1.3}{The \underline{RAM Model of Computation} is a highly simplified model of computation that we use to compute the efficiency of an algorithm. Its memory consists of an infinite array which it uses to carry out instructions sequentially. All instructions take unit time, meaning the running time is the number of instructions executed.}

\separate{5}\definition{1.1.4}{\underline{Pseudocode} is a generic way of describing how an algorithm will function without choosing any specific programming language. It does not have any very strict rules, but follows the following generic format:}

\algorithmBlock{Example pseudocode}{Input conditions}{The algorithm's output}
{
    \State \(x \gets 1\)
    \State \(y \gets 2\)

    \If {\(y > x\)}
        \State \Return y
    \Else
        \State \Return x
    \EndIf
}

\noindent\definition{1.1.5}{A \underline{variable} is used to store some value, for example, in the above code, \(x\) and \(y\) are variables. Some basic types include integer, float, char, and string.}

\subsection{Recursion and Backtracking}

\definition{1.2.1}{A \underline{recursive algorithm} is an algorithm which calls itself to do a part of its work. It should have a base case (where the final result is returned), must change its state to move towards the base case, and must call itself recursively.}

\separate{5}\definition{1.2.2}{A \underline{backtracking algorithm} tries to construct a solution to a computational problem incrementally, one small piece at a time. Whenever the algorithm needs to decide between multiple alternatives to the next component of the solution, it recursively evaluates every alternative and then chooses the best one.} Definition from \hyperref[cit:1.algorithms]{(Erickson, 2019, p.71)}.

\separate{5}\definition{1.2.3}{\underline{Memoisation} is the process of storing the result of a computation so that it can be subsequently retrieved without repeating the computation.}

\newpage\section{Basic Data Structures}

\subsection{Arrays}

\definition{2.1.1}{An \underline{array} is a sequence of elements that is usually denoted as \(A[1], \; A[2], \; ..., \; A[n]\) or as \(A[1,...,n]\).} The elements are stored in consecutive memory cells and are all of the same data type. 

\subsection{Lists}

Lists can be separated into three different categories, depending on how they store data. 

\separate{5}\definition{2.2.1}{A \underline{singly-linked list} consists of nodes, each storing an element and a pointer to the subsequent node in the list. This means that the nodes do not have to be stored consecutively.}
\begin{itemize}
    \item \definition{2.2.2}{The first node in a singly-linked list is called the \underline{head}.}
    \item \definition{2.2.3}{The final node in a singly-linked list is called the \underline{tail} and its pointer points to null.}
\end{itemize}
\definition{2.2.4}{A \underline{doubly-linked list} consists of nodes with a value and two pointers, pointing to both the previous and subsequent nodes in the list. The list also must keep a size counter to track the number of nodes in the list.}
\begin{itemize}
    \item \definition{2.2.5}{At either end of the doubly-linked list is a \underline{sentinel node}. These nodes either have a null next or previous pointer and do not store a value. An empty list has two sentinel nodes pointing at each other.}
\end{itemize}
\definition{2.2.6}{A \underline{circularly-linked} list uses the same nodes as a singly-linked list but does not have a head or tail node. Instead, the final node points to the start of the list, creating a circle.}
\begin{itemize}
    \item \definition{2.2.7}{One node in the circularly-linked list is designated as the \underline{cursor} and is used as a starting point when traversing the list.}
\end{itemize}

\subsection{Stacks}

\definition{2.3.1}{A \underline{stack} is a collection of objects that are inserted and removed according to the last-in-first-out (LIFO) principle.} They support the following methods:
\begin{itemize}
    \item \textbf{Compulsory Methods:}
    \begin{itemize}
        \item \textbf{Push}: \(push(element)\) inserts \(element\) at the top of the stack.
        \item \textbf{Pop}: \(pop()\) removes and returns the top element from the stack. An error occurs if the stack is empty.
    \end{itemize}
    \item \textbf{Optional Methods:}
    \begin{itemize}
        \item \textbf{Size}: \(size()\) returns the number of elements in the stack.
        \item \textbf{Is Empty}: \(isEmpty()\) returns true if the stack is empty, and false if not.
        \item \textbf{Top}: \(top()\) returns the top item from the stack without removing it. Throws an error if the stack is empty.
    \end{itemize}
\end{itemize}
Stacks are generally implemented using an array to store the items of the stack; however this can lead to memory wastage when the stack is not full, and conversely it will throw errors if the stack is full.

\newpage\subsection{Queues}

\definition{2.4.1}{A \underline{queue} is a collection of objects that are inserted and removed according to the first-in-first-out (FIFO) principle.} They support the following methods:
\begin{itemize}
    \item \textbf{Compulsory Methods:}
    \begin{itemize}
        \item \textbf{Enqueue}: \(enqueue(element)\) inserts \(element\) at the end of the queue.
        \item \textbf{Dequeue}: \(dequeue()\) removes the element from the front of the queue. An error occurs if the queue is empty.
    \end{itemize}
    \item \textbf{Optional Methods:}
    \begin{itemize}
        \item \textbf{Size}: \(size()\) returns the number of elements in the queue.
        \item \textbf{Is Empty}: \(isEmpty()\) returns true if the queue is empty, and false if not.
        \item \textbf{Front}: \(front()\) returns the element from the front of the queue without removing it. Throws an error if the queue is empty.
    \end{itemize}
\end{itemize}

\subsection{Hash Tables}

\definition{2.5.1}{A \underline{hash table} consists of a bucket array and a hash function.}

\separate{5}\definition{2.5.2}{A \underline{bucket array} for a hash table is an array \(A\) of size \(N\), where each cell of \(A\) is thought of as a bucket storing a collection of key-value pairs.} 

\separate{5}\definition{2.5.3}{The size \(N\) of the array is called the \underline{capacity} of the hash table.}

\separate{5}\definition{2.5.4}{A \underline{hash function} \(h\) is a function mapping each key \(k\) to an integer in the range \([0,N-1]\), where \(N\) is the capacity of the hash table. This provides an index \(h(k)\) of a location in \(A\) where the key-value pair \((k,v)\) is stored. Hash functions are usually composed of the hash code (mapping keys to integers) and a compression function (integers to indices).}

\separate{5}\definition{2.5.5}{A \underline{collision} occurs when two keys map to the same hash value, and consequently want to be stored in the same bucket. These can be dealt with using different methods:}
\begin{itemize}
    \item \definition{2.5.6}{\underline{Separate chaining} converts each bucket to a list of key-value pairs. This will result in lists of average length \(\frac{N}{M}\) where \(N\) is the amount of data, and \(M\) is the capacity of the hash table.}
    \item \definition{2.5.7}{If a collision occurs at \(h(k)=i\), \underline{linear probing} tries to insert again at \(A[(i+1)\%N]\), then if that is also full, tries at \(A[(i+2)\%N]\) etc. until the key-value pair is inserted.}
    \item \definition{2.5.8}{\underline{Quadratic probing} iteratively tries the buckets \(A[(i+f(j)) \; \% \; N]\) for \(j=0,1,2,...\) where \(f(j)=j^2\) until finding an empty bucket. This avoids clustering patterns that occur with linear probing. However, it creates secondary clustering. This strategy may not find an empty bucket in $A$ even if one exists.}
    \item \definition{2.5.9}{Using \underline{double hashing}, we choose a secondary hash function $h'$, and if $h$ maps some key $k$ to a bucket $A[i]$, with $i=h(k)$, that is already occupied, then we iteratively try the buckets $A[(i+f(j)) \; \% \; N]$ for $j=0,1,2,...$, where $f(j)=j \cdot h'(k)$. A common choice of secondary hash function is $h'(k)=q-(k \; \% \; q)$, for some prime number $q<N$. The secondary hash function should not evaluate to 0. }
\end{itemize}
\definition{2.5.10}{Deletions from the hash table must not hinder future searches. To ensure future searches are not hindered, we use \underline{tombstones}. These are markers that are left in a bucket after a deletion. If a tombstone is encountered when searching along a probe sequence to find a key, we know to continue. The use of tombstones lengthens the average probe sequence distance. Two possible remedies are:}
\begin{itemize}
    \item \definition{2.5.11}{Using \underline{local reorganisation}, after deleting a key, we continue to follow the probe sequence of that key and move records into the vacated bucket.} 
    \item Or, we can periodically rehash the table by reinserting all records into a new hash table. If you have a record of which keys are accessed most, these can be placed where they can be found most easily. 
\end{itemize}

\newpage\part{Asymptotic Notation and Sorting}\label{part:2.ans}

\section{Asymptotic Notation}

\subsection{Types of Complexity}

\definition{3.1.1}{The \underline{time complexity} of an algorithm can be expressed as the number of basic operations executed by an algorithm when computing for a set input size.}

\separate{5}\definition{3.1.2}{The \underline{space complexity} of an algorithm is the amount of memory required to carry out the algorithm on an input of a set size.}

\separate{5}\definition{3.1.3}{The \underline{worst-case time complexity} of an algorithm is the largest number of basic operations possible for an input of a set size.}

\subsection{General Notation}

\definition{3.2.1}{Generally, \underline{Big-O Notation} denotes that, if \(f = O(g)\), then \(f\) will grow no faster than \(g\) (the upper-bound for the growth rate, or \(f' \leq g'\)).} Formally, let \(f\) be a real or complex function, and let \(g\) be a real function such that
\begin{equation}
    f(x) = O(g(x)) \quad\text{as } x \rightarrow \infty
\end{equation}
meaning \(f(x)\) is Big-O \(g(x)\) if the absolute value of \(f(x)\) is at most a positive constant multiple of the absolute value of \(g(x)\) for a suitably large value of \(x\):
\begin{equation}
    |f(x)| \leq C \cdot |g(x)| \quad x \geq x_0, \; x_0 \in \R, \; C \in \N
\end{equation}

\separate{5}\definition{3.2.2}{Similarly, \underline{Big-Omega Notation} denotes that, if \(f=\Omega(g)\), then \(f\) will grow at least as fast as \(g\) (the lower-bound for the growth rate, or \(f' \geq g'\)).} Formally, let \(f\) and \(g\) be functions \(\R \rightarrow \R\). We say that 
\begin{equation}
    f(x) = \Omega(g(x))
\end{equation}
if, for a suitably large value of \(x\), there exists a constant \(C\) such that
\begin{equation}
    |f(x)| \geq C \cdot |g(x)| \quad x \geq x_0, \; x_0 \in \R, \; C \in \N
\end{equation}
Therefore, \(f(x)\) is \(\Omega(g(x))\) if \(g(x)\) is \(O(f(x))\).

\separate{5}\definition{3.2.3}{\underline{Big-Theta Notation} denotes the exact growth rate of a function. This means that, for some real-valued functions \(f\) and \(g\),}
\begin{equation}
    f(x) = O(g(x)) \quad\text{and}\quad f(x)=\Omega(g(x))
\end{equation}

\separate{5}\definition{3.2.4}{\underline{Little-o Notation} can be used for disregarding \say{small-order} terms. Generally, \(f(x) = o(g(x)\) if \(f'(x) < g'(x)\). It is formally defined for the real-valued functions \(f\) and \(g\) such that \(f(x) = o(g(x))\) when}
\begin{equation}
    \lim_{x\rightarrow\infty} \frac{f(x)}{g(x)} = 0
\end{equation}
When the limit is removed, this becomes
\begin{equation}
    o(g) = \{f: \N \rightarrow \N \;|\; \forall C > 0 \;\exists k > 0 : C \cdot f(n) < g(n) \;\forall n \geq k \} 
\end{equation}
Therefore, \(f(x) = o(g(x) \implies f(x) = O(g(x))\).

\newpage\noindent\definition{3.2.5}{Similar to Little-o Notation, \underline{Little-omega Notation} is defined such that \(f=\omega(g)\) if \(f' > g'\),}
and is defined formally as
\begin{equation}
    \omega(g) = \{f: \N \rightarrow \N \;|\; \forall C > 0 \;\exists k > 0 : f(n) > C \cdot g(n) \;\forall n \geq k \} 
\end{equation}

\separate{5}\definition{3.2.6}{A function is designated as a \underline{sublinear function} if it grows slower than any linear function. Using Little-o Notation, \(f(x)\) is sublinear if \(f(x) = o(x)\), meaning}
\begin{equation}
    \lim_{x \rightarrow \infty} \frac{f(x)}{x} = 0
\end{equation}

\section{Sorting}

\subsection{Common Algorithms}

\definition{4.1.1}{An \underline{insertion sort} works by taking the \(i^{th}\) element from a list and inserting it into an already sorted sequence of elements \([1,\;j-1]\).}
This has a best-case runtime of \((n-1)\) and a worst-case of \(\frac{n(n-1)}{2}\), giving a Big-O of \(O(n^2)\).

\separate{5}\definition{4.1.2}{A \underline{selection sort} takes the \(i^{th}\) smallest elements and places them in order at the start of the list until all elements have been sorted.}
Its time complexity can be calculated by doing the following:
\begin{equation}
    \sum^{n-1}_{i=1} \sum^n_{j=i+1} 1 
    = \sum^{n-1}_{i=1} (n-i) 
    = \left( \sum^{n-1}_{i=1} n \right) - \left( \sum^{n-1}_{i=1} i \right) 
    = n(n-1) - \frac{n(n-1)}{2}
    = O(n^2)
\end{equation}

\separate{5}\definition{4.1.3}{A \underline{bubble sort} takes the \(i^{th}\) largest elements and places them in order at the end of the list until all elements have been sorted.}
Its time complexity can be calculated by doing the following:
\begin{equation}
    \sum^{n-1}_{i=1} (n-i) = O(n^2)
\end{equation}

\separate{5}\definition{4.1.4}{A \underline{merge sort} takes a list of elements and recursively splits it into two even parts until each sublist contains one element. Once this occurs, the sublists are merged back together in the correct order to form the sorted list. It does this with a variant of an insertion sort.}
It has a time complexity of \(O(n \, log \, n)\).

\separate{5}\definition{4.1.5}{A \underline{quick sort} takes the first, middle, and final element of a list and selects one as a pivot point. The remaining elements are then placed in sublists depending on whether they are smaller/larger than the pivot.}
This gives a best-case time complexity of \(o(n \, log \, n)\) and a worst-case of \(O(n^2)\).

\newpage\subsection{Recursion for Sorting and Solving Recurrences}

To analyse recursive algorithms, first of all note that most of them are actually hybrids between iterative and strictly recursive. For example, merge sort splits the input of size \(n\) into two halves of equal size \(\frac{n}{2}\), independently recurses in each half, then merges the resulting sorted sequence.

\begin{equation}
    T(n) \leq 
    \begin{cases}
        d   \qquad \text{if } n \leq c \text{ for constants } c,\,d > 0 \\
        2 \cdot T(\frac{n}{2}) + a \cdot n  \qquad \text{otherwise}
    \end{cases}
\end{equation}
\definition{4.2.1}{The \underline{Master Theorem} can be used to quickly solve recurrences of form}
\begin{equation}
    T(n) \leq a T\left(\frac{n}{2}\right) + f(n) \quad\quad a \geq 1, \, b > 1
\end{equation}
The theorem has three cases:
\begin{enumerate}
    \item \(f(n) = O(n^{log_b a - \epsilon}), \; \epsilon > 0 \implies T(n) = \Theta(n^{log_b a})\)
    \item \(f(n) = \Theta(n^{log_b a} \cdot (log \, n)^k) \implies T(n) = \Theta(n^{log_b a} \cdot (log \, n)^{k+1})\)
    \item \(f(n) = O(n^{log_b a + \epsilon}), \; \epsilon > 0\) and \(a f(\frac{n}{b}) \leq c f(n), \; c < 1 \implies T(n) = \Theta(f(n))\)
\end{enumerate}

\newpage\part{Selection and Data Structures}\label{part:3.sds}

\section{Algorithms for Sorting, Searching, and Selection}

\subsection{Sorting}

\definition{5.1.1}{A \underline{bucket sort} algorithm creates buckets with labels in the interval \([1,k]\) where \(k\) is the largest value in the list. Each element is then placed in its corresponding bucket before the buckets are outputted in ascending order, creating a sorted list.}
This gives a worst-case time complexity of \(O(n+k)\). If \(k\) is small relative to \(n\), this gives a best-case time complexity of \(o(n \, log \, n)\). 

\separate{5}\definition{5.1.2}{A \underline{bucket sort} recursively sorts a list of items by a given digit of each element until the list is sorted (from right-most digit to left-most). This algorithm requires as many buckets as the base of the number system being used (e.g. binary requires two buckets).}
It has a time complexity of \(O(n \, log_d \, k)\) for a list of values in \([1,k-1]\) in base \(d\).

\subsection{Searching}

\definition{5.2.1}{A \underline{binary search} assumes that the list being searched is sorted. It recursively halves the search space on each pass until the desired element is found. This produces the recurrence}
\begin{equation}
    T(n) = T\left(\frac{n}{2}\right) + O(1) = O(log \, n)
\end{equation}

\subsection{Selection}

Selection algorithms are used to, for example, find the \(i^{th}\) smallest element in an unsorted list.

\separate{5}\definition{5.3.1}{The \underline{quick select} algorithm works similarly to the quick sort, using pivots to split the search space and reduce execution times.}
In the worst-case, the algorithm is \(O(n^2)\), but on average executes in \(O(n)\) time.

\separate{5}\definition{5.3.2}{The \underline{median-of-medians} algorithm is doubly recursive and works by calculating the median of the list and using it as a pivot.} This gives the recurrence
\begin{equation}
    T(n) = T\left(\frac{n}{5}\right) + T\left(\frac{7n}{10} + 3\right) + n = O(n)
\end{equation}

\section{Trees}

\subsection{Binary Search Trees}

\definition{6.1.1}{A \underline{rooted binary tree} is a tree where each node has a single parent node and at most two children nodes. The binary tree is rooted if there is one node with all other nodes as descendants.}

\separate{5}\textbf{Tree Properties}:
\begin{itemize}
    \item \definition{6.1.2}{A tree's \underline{size} is the total number of nodes in the tree.}
    \item \definition{6.1.3}{A node's \underline{depth} is its distance from the root node.}
\end{itemize}

\separate{5}\definition{6.1.4}{A binary tree becomes a \underline{binary search tree} if, for all nodes, all left-descendant nodes are smaller than it, and all right-descendant nodes are greater than it.}

\newpage\noindent\theorem{6.1.1}{Let \(T\) be a rooted binary tree of height \(h\). Then,
\begin{itemize}
    \item \(T\) has at most \(\left(2^{h+1}-1\right)\) nodes.
    \item \(T\) has at most \(2^h\) leaves.
\end{itemize}}
{The maximum number of nodes is in a complete tree of height \(h\), meaning all levels are full,
\begin{equation}
    1 + 2 + 2^2 + 2^3 + ... + 2^h 
    = \sum^h_{n=0} 2^n 
    = 2^{h+1} - 1
\end{equation}
The second statement follows by induction on \(h\). Case \(h=0\) is obvious. We then consider the left and right subtrees of the root of \(T\).
Then each of them has height \(\leq (h-1)\), and (by induction hypothesis) \(\leq 2^{h-1}\) leaves. Then \(T\) has at most \(2 \times 2^{h-1} = 2^h\) leaves.}

\separate{5}Trees can be traversed in one of three ways, these are:
\begin{itemize}
    \item \definition{6.1.5}{\underline{Preorder traversal} where the root is visited, followed by the left and right subtrees.}
    \item \definition{6.1.6}{\underline{Inorder traversal} where the root is visited in between the left and right subtrees.}
    \item \definition{6.1.7}{\underline{Postorder traversal} where the root is visited after both the left and right subtrees.}
\end{itemize}

\subsection{Balanced Binary Search Trees}

\definition{6.2.1}{\underline{Balanced binary search trees} are binary search trees which balance themselves using rotation operations to minimise height differences between subtrees.}

\image{Images/Balanced BSTs - Left and right rotation examples.png}{The basic rotation operations.}{0.75}

\newpage\subsection{AVL Trees}

\definition{6.3.1}{An \underline{AVL tree} is a self-balancing binary search tree with the height-balance property (for each node, the height of its children differ by at most one).}

\separate{5}\theorem{6.3.1}{The height of an AVL tree storing \(n\) nodes is \(O(log\,n)\)}
{Let us bound \(n(h)\) as the minimum number of internal nodes of an AVL tree of height \(h\). We can clearly see that \(n(1)=1\) and \(n(2)=2\). For \(n > 2\), an AVL tree of height \(h\) contains the root node, one AVL subtree of height \(n(h-1)\), and another of height \(n(h-2)\). That is,
\begin{equation}
    n(h) = n(h-1) + n(h-2)
\end{equation}
Knowing that \(n(h-1) > n(h-2)\), we get \(n(h) > 2n(h-2)\) so, by induction,
\begin{equation}
    n(h) > 2^in(h-2i)
\end{equation}
Taking \(i\) to be \(\frac{h}{2}-1\), so that \(h-2i = 1\) or 2, on solving we get,
\begin{equation}
    n(h) > 2^{\frac{h}{2}-1}
\end{equation}
Taking logarithms, \(h < 2log(n(h)) + 2\). Since \(log(n) > log(n(h))\), the height of an AVL tree is \(O(log\,n)\).}

\separate{5}\definition{6.3.2}{The algorithm for \underline{AVL insertion} automatically rebalances the tree after inserting a new element to uphold the height-balance property. For example, for a new node \(n\),
\begin{itemize}
    \item If it was inserted into the left subtree of a left child, the algorithm performs a single right rotation.
    \item If it was inserted into the right subtree of a right child, the algorithm performs a single left rotation.
    \item If it was inserted into the right subtree of a left child, the algorithm performs a left rotation followed by a right rotation.
    \item If it was inserted into the left subtree of a right child, the algorithm performs a right rotation followed by a left rotation.
\end{itemize}}

\section{Heaps}

\subsection{Heap Properties}

\definition{7.1.1}{A \underline{priority queue} is a collection of prioritised elements that allows arbitrary element insertion and removal having the first priority (highest or lowest).}

\separate{5}\definition{7.1.2}{All \underline{heaps} are complete binary trees. The first item in a heap is the root, the second is the left child, and the third is the right child. The next nodes always fill the next levels left-to-right.}

\separate{5}\definition{7.1.3}{A \underline{max heap} stores values where, for every node except the root, the value is at most that of its parent. In any subtree, no values are larger than the value stored in the subtree root.}

\separate{5}\definition{7.1.4}{A \underline{min heap} stores values where, for every node except the root, the value is at least that of its parent. In any subtree, no values are smaller than the value stored in the subtree root.}

\separate{5}\definition{7.1.5}{Assuming that we start counting at one for the root, the \underline{heap indices} can be calculated such that, for a node at index \(i\), the left child is at \(2i\), the right child is at \(2i+1\), and the parent is at \(\left\lfloor\frac{i}{2}\right\rfloor\)}.

\newpage\subsection{Heap Operations}

\definition{7.2.1}{\underline{\hyperref[alg:Heapify]{Heapify}} takes an incorrectly ordered heap and reorders it into an acceptable order. It has a runtime of \(O(log\,n)\).}

\algorithmBlock{Heapify}
{\(A\), the array to turn into a heap, \(v\), the index of the largest value, \(i\), the highest index.}
{A correctly ordered heap.}
{
    \State largest \(\gets v\)
    \State leftChild \(\gets 2v\)
    \State rightChild \(\gets 2v+1\)

    \separate{5}\If {leftChild \(\leq i\) and \(A[\text{leftChild}] > A[\text{largest}]\)}
        \State largest \(\gets\) leftChild
    \EndIf

    \If {rightChild \(\leq i\) and \(A[\text{rightChild}] > A[\text{largest}]\)}
        \State largest \(\gets\) rightChild
    \EndIf

    \separate{5}\If {largest \(\neq v\)}
        \State Swap(A[v], A[largest])
        \State Heapify(A, largest, i)
    \EndIf
}

\noindent\definition{7.2.2}{\underline{Insertion} adds the new element to the bottom-left of the heap before running \hyperref[alg:Heapify]{Heapify}.}

\separate{5}\definition{7.2.3}{\underline{Build heap} takes an array of numbers and converts it into a heap by running \hyperref[alg:Heapify]{Heapify} iteratively, starting \(v\) as the size of the array and looping down to one.}

\separate{5}\theorem{7.2.4}{The build heap algorithm runs in \(O(n)\) time.}
{Let \(T(n)\) be the execution time for a heap of size \(n\).
    \begin{equation}
        T(n) = \sum^{log\,n}_{h=1} (\# \text{ of nodes at height h}) \cdot O(h)
    \end{equation}
    Given the properties of binary trees, 
    \begin{equation}
        T(n) \leq \sum^{log\,n}_{h=1} \frac{n}{2^h} \cdot O(h) = O\left( n \cdot \sum^{log\,n}_{h=1} \frac{h}{2^h} \right)
    \end{equation}
    This is well known for \(x \in \{0,1\}\),
    \begin{equation}
        \sum^\infty_{i=0} i \cdot x^i = \frac{x}{(1-x)^2} \quad\quad\quad \sum^{log \, n}_{h=1} \frac{h}{2^h} \leq \sum^\infty_{h=0}h \cdot \left(\frac{1}{2}\right)^h = 2
    \end{equation}
    and hence,
    \begin{equation}
        T(n) = O \left(n \cdot \sum^{log \, n}_{h=1} \frac{h}{2^h}\right) = O(n)
    \end{equation}
}

\newpage\section{Lower Bounds for Sorting and Selection}

\theorem{8.0.1}{Any comparison-based sorting algorithm requires \(\Omega(n\,log\,n)\) comparisons in the worst case.}
{It is sufficient to determine the minimum height of a decision tree in which each permutation appears as a leaf. 
We can consider a decision tree of height \(h\) with \(l\) leaves corresponding to a comparison sort on \(n\) elements. Each of the \(n!\) permutations of the input appears as some leaf, \(l \geq n!\). Given that a binary tree of height \(h\) has at most \(2^h\) leaves, we get that
\begin{equation}
    n! \leq l \leq 2^h
\end{equation}
Therefore, \(n! \leq 2^h\), and after taking logs of both sides,
\begin{equation}
    h \geq log(n!) \equiv \Omega(n\,log\,n)
\end{equation}}

\separate{5}\definition{8.0.1}{An \underline{adversary} algorithm is designed to intercept and limit access to the input by the original algorithm such that the time taken to make a decision is maximised. However, the adversary does not have knowledge of how the algorithm works, so must be general enough to work for any algorithm provided.}

\separate{5}\theorem{8.0.2}{Any comparison-based algorithm attempting to find the maximum element in an unsorted array requires at least \((n-1)\) comparisons in the worst case.}
{After at most \((n-2)\) comparisons no more than two elements can have not lost a comparison. An adversary can make either of them max and be consistent, but there is not enough information for the original algorithm to be consistent. Hence, the algorithm needs to make one more comparison, making the minimum number of comparisons \((n-1)\).}
\textit{Alternative }\proof{alt:thm:8.0.2}{Taking this into account, we can design a strategy for the adversary. The only relevant information for the algorithm is how many elements have lost at least one comparison. As soon as \((n-1)\) elements have lost, the algorithm can decide that the only undefeated element must be the maximum. 

Assuming each index \(i\) has a status (\(N\) for never lost, and \(L\) for lost), initially all indices have \(N\). After each comparison there can be a status change \(N\rightarrow L\). Therefore, the adversary must delay changes \(N\rightarrow L\) as much as possible. 

Given this information, each comparison will either update one status from \(N\) to \(L\), or will update neither, providing the algorithm with at most one bit of new information. In total, the algorithm needs \((n-1)\) bits of information to make a decision, and hence at least \((n-1)\) comparisons are necessary.}

\newpage\part{Graph Theory and Algorithms}\label{part:4.gta}

\section{The Basics of Graph Theory}

\subsection{Types of Graphs}

\definition{9.1.1}{A \underline{graph} \(G\) is a pair \((V(G), E(G))\) of sets where \(V(G)\) is a non-empty set of vertices making up the graph and \(E(G)\) is a set of unordered pairs, \(\{u,v\}, \;u,v\in G\).}

\separate{5}\definition{9.1.2}{A \underline{directed graph} is a graph where the edge set \(E(G)\) instead contains ordered pairs, indicating the direction of travel along each edge.}

\separate{5}\definition{9.1.3}{A \underline{multi-graph} allows for multiple edges to exist between vertices, meaning the edge set may contain duplicates.}

\separate{5}\definition{9.1.4}{A \underline{pseduo-graph} allows loops, meaning pairs in the edge set can contain two instances of the same vertex.}

\separate{5}\definition{9.1.5}{A \underline{vertex/edge-weighted graph} assigns a weight for each vertex or edge.}

\separate{5}\definition{9.1.6}{A \underline{complete bipartite graph} is a graph such that any subgraph is a \underline{bipartite graph}, meaning we can separate its vertex set into two separate sets where each edge has an endpoint in both sets.}

\separate{5}\definition{9.1.7}{A \underline{complete graph} contains all possible edges between vertices.}

\separate{5}\definition{9.1.8}{The \underline{hypercube, n-dimensional cube, or n-cube} is the graph with
\begin{equation}
    V = \{(e_1,...,e_n) | e_i \in \{0,1\} (i=1,...,n)\}
\end{equation}
in which two vertices are neighbours iff the corresponding rows differ in exactly one entry.}

\subsection{Other Terminology}

Let \(G\) be a graph and \(uv\) be an edge in it. Then,
\begin{itemize}
    \item \definition{9.2.1}{\(u\) and \(v\) are called \underline{endpoints} of the edge \(uv\).}
    \item \definition{9.2.2}{\(u\) and \(v\) are called \underline{neighbours} or \underline{adjacent vertices}.}
    \item \definition{9.2.3}{\(uv\) is said to be incident to \(u\) and \(v\).}
    \item \definition{9.2.4}{If \(uw\) is also an edge where \(v\neq w\), then \(uv\) and \(uw\) are \underline{adjacent edges}.}
\end{itemize}

\separate{5}\definition{9.2.5}{The \underline{neighbourhood} of a vertex \(v \in V\), denoted by \(N(v)\), is the set of neighbours of \(v\).}

\separate{5}\definition{9.2.6}{The \underline{degree} of a vertex \(v \in V\), \(deg(v)\), is the size of its neighbourhood, \(|N(v)|\).}

\separate{5}\definition{9.2.7}{\underline{\(\delta(G) \text{ and } \Delta(G)\)} give the vertices with the smallest and largest degrees, respectively.}

\separate{5}\definition{9.2.8}{A vertex with degree zero is an \underline{isolated vertex}.}

\separate{5}\definition{9.2.9}{A vertex with degree one is an \underline{end vertex} or a \underline{pendant vertex}.}

\separate{5}\definition{9.2.10}{A \underline{subgraph} \(G'=(V',E')\) of \(G=(V,E)\) is a graph with \(V' \subseteq V\) and \(E' \subseteq E\).}
\begin{itemize}
    \item \definition{9.2.11}{A subgraph is \underline{proper} if \(G'\neq G\).}
    \item \definition{9.2.12}{A subgraph is \underline{spanning} if \(V' = V\).}
    \item \definition{9.2.13}{An \underline{induced subgraph} is a subgraph where \(E'\) contains all edges of \(E\) between vertices of \(V'\). It is obtained by removing from \(G\) all vertices of \(G \backslash G'\).}
\end{itemize}

\subsection{Basic Theorems}

\separate{5}\theorem{9.3.1}{Known as the \underline{Handshaking Lemma}, if \(G=(V,E)\) is a graph, then \(\sum_{v\in V} deg(v) = 2|E|\).}
{Every edge has two endpoints and contributes one to each of their degrees, so contributes two to the sum of the degrees of all vertices of \(V\).}

\separate{5}\theorem{9.3.2}{In every undirected graph, the number of vertices of odd degree is even.}
{Let \(G=(V,E)\) and \(V_{odd} = \{v:deg(v) \text{ is odd}\}\) and \(V_{even} = \{v:deg(v) \text{ is even}\}\). Clearly \(\sum_{v\in V_{even}} deg(v)\) is even. By the \hyperref[thm:9.3.1]{Handshaking Lemma},
\begin{equation}
    \sum_{v \in V_{odd}} deg(v) = 2 \cdot |E| - \sum_{v \in V_{even}}
\end{equation}
it follows that \(\sum_{v \in V_{odd}} deg(v)\) is even too. Now, if we have an odd number of vertices of odd degree, then there is a contradiction, meaning there must be an even number of vertices with an odd degree.} 

\separate{5}\theorem{9.3.3}{All n-dimensional cubes are bipartite.}{Let \(K_n\) be an n-dimensional hypercube. Create two vertex sets, \(V_1\) and \(V_2\), containing all the vertices of \(K_n\) with even and odd numbers of ones, respectively. Due to the definition of a hypercube, it is guaranteed for each edge to have an endpoint in both of the resulting vertex sets. Therefore, the n-cube must be bipartite.}

\section{Paths, Cycles, and Connectivity}

\subsection{Walks, Paths, and Cycles}

\definition{10.1.1}{A \underline{walk} is a sequence of edges from a graph \(G\).}

\separate{5}\definition{10.1.2}{A \underline{path} in a graph is a subgraph which is isomorphic to the graph \(P_k\) for some integer \(k \geq 1\).}

\separate{5}\definition{10.1.3}{A \underline{circuit} is a walk where the starting node is also the ending node.}

\separate{5}\definition{10.1.4}{A \underline{cycle} is a subgraph isomorphic to the graph \(C_k\) where \(k \geq 3\).}

\separate{5}\definition{10.1.5}{The \underline{length} of a path is the number of edges within it.}

\separate{5}\definition{10.1.6}{The \underline{distance} between two nodes, \(dist(u,v)\), is the length of the shortest path between the two vertices, and is defined as \(\infty\) if no path exists.}

\separate{5}\definition{10.1.7}{The \underline{diameter} of a graph is the largest possible distance between two vertices in the graph.}

\subsection{Connectivity}

\definition{10.2.1}{A graph is \underline{connected} if, between every pair of vertices, there exists at least one path.}

\separate{5}\definition{10.2.2}{A \underline{connected component} of a graph is a maximally connected subgraph.}

\newpage\noindent\theorem{10.2.1}{Every graph contains at least \(|V|-|E|\) connected components.}
{We can prove this by induction over \(m=|E|\). When \(m=0\) there are no edges, only isolated vertices, meaning there must be \(|V|\) connected components.

Let \(G=(V,E)\) with \(|E|=m+1\). Consider an arbitrary \(e\in E\), and define \(E'=E \backslash \{e\}\). By the inductive hypothesis, \(G'=(V,E')\) has at least \(|V|-|E'|=|V|-m\) connected components. This produces two cases: either, with \(e\), the number of connected components does not change; or, with \(e\), the number of connected components decreases by one. In both cases, \(G\) has at least \(|V|-m-1=|V|-|E|\) connected components, proving the initial hypothesis.}

\separate{5}\definition{10.2.3}{A directed graph is \underline{weakly connected} if the graph obtained from \(G\) by forgetting directions is connected.}

\separate{5}\definition{10.2.4}{A directed graph is \underline{strongly connected} if any two distinct vertices are connected by paths in both directions.}

\separate{5}\definition{10.2.5}{A \underline{strongly connected component} of a directional graph is a maximal strongly connected subgraph.}

\subsection{Special Circuits and Cycles in Graphs}

\definition{10.3.1}{A \underline{Eulerian circuit} is a graph where we can start and end at the same node while traversing each edge exactly once.}

\separate{5}\definition{10.3.2}{A \underline{Hamiltonian cycle} is a graph where we can start and end on the same node while visiting each node exactly once.}

\separate{5}\theorem{10.3.1}{A connected graph with at least two vertices has an Eulerian circuit iff each of its vertices has even degree.}
{This is proved in \hyperref[cit:2.eulerianCircuits]{(Rosen, 2003, pp.634-636)}.}

\section{Trees and Isomorphism}

\subsection{Trees}

\definition{11.1.1}{A \underline{forest} is an acyclic graph.}

\separate{5}\definition{11.1.2}{A \underline{tree} is a connected forest.}

\separate{5}\theorem{11.1.1}{Every connected graph contains a spanning tree.}{Let \(G\) be a connected graph. If \(G\) doesn't contain cycles then it is a tree and therefore is its own spanning tree. Conversely, if the graph does contain a cycle, then we can remove one edge from the graph, breaking the cycle and forming a spanning tree. The resulting graph will still be connected as cycles will contain at least one vertex of degree two, meaning removing an edge won't produce an isolated vertex.}

\newpage\noindent\definition{11.1.3}{A \underline{leaf} is a vertex of degree one.}

\separate{5}\theorem{11.1.2}{Every tree on at least two vertices contains a leaf.}{By contradiction, we assume that every vertex has degree zero or two. If a vertex has degree zero, then the graph is not connected, and therefore cannot be a tree. Similarly, if every vertex has degree two then we can traverse through the tree by travelling to each vertex's neighbours until we reach a node that has already been visited, proving a cycle exists, making the graph not a tree. Therefore, trees with at least two vertices must contain at least one leaf.}
\textit{Alternative }\proof{alt:thm:11.1.2}{(When every vertex has a degree of at least two). Consider a longest path \(P\) and an end vertex \(v\) of \(P\). All neighbours of \(v\) are on \(P\) as the tree must be connected, therefore if \(deg(v) \geq 2\), then there is a cycle, meaning a tree cannot exist. This holds for both endpoints of the path, proving there must be at least one vertex of degree one for the graph to be a tree.}

\separate{5}\theorem{11.1.3}{A connected graph on \(n\) vertices is a tree iff it has \((n-1)\) edges.}
{For small values of \(n\), the theorem holds, e.g., a tree with one vertex has no edges, and a tree with two vertices has one edge. Suppose each tree on \((n-1)\) vertices has \((n-2)\) edges. Take a tree on \(n \geq 3\) vertices, \(T\) contains a leaf \(v\). Consider the graph \(T-v\), it has one vertex less and one edge less than \(T\). \(T-v\) is still acyclic and such is a tree with \(n-1\) vertices, by the inductive hypothesis it has \((n-2)\) edges. \(T\) has one edge more, so \((n-1)\) edges.}

\separate{5}\theorem{11.1.4}{Let \(T\) be a tree and \(u,v \in V(T)\) with \(u \neq v\). Then there is a unique path in \(T\) between \(u\) and \(v\).}
{By contradiction, suppose there are two paths \(P\) and \(Q\) in \(T\) between \(u\) and \(v\). Let \(x\) and \(y\) in \(V(T)\) be distinct and chosen in such a way that \(x\) and \(y\) are on \(P\) and \(Q\), but between \(x\) and \(y\) the vertices of \(P\) and \(Q\) are disjoint. Then the segments of \(P\) and \(Q\) form a cycle, contradicting that \(T\) is a tree.}

\separate{5}\definition{11.1.4}{A \underline{rooted tree} is a tree in which one vertex is fixed as the root (and every other vertex is directed away from the root).}

\separate{5}Let \(v\) be a vertex in a rooted tree \(T\).
\begin{itemize}
    \item \definition{11.1.5}{The neighbours of \(v\) in the next level are called the \underline{children} of \(v\).}
    \item \definition{11.1.6}{The (unique) neighbour of \(v\) in the previous level is called the \underline{parent} of \(v\).}
    \item \definition{11.1.7}{If \(v\) has children, then it is an \underline{internal} vertex.}
\end{itemize} 

\separate{5}\theorem{11.1.5}{A tree with at least two vertices, \(n_3\) of which have degree at least three, has at least \((n_3+2)\) leaves.}
{Let \(T\) be a tree on \(n \geq 2\) vertices. We use induction on \(n\). Let \(l(t)\) denote the number of leaves in \(T\), and \(n_3(T)\) denote the number of vertices of degree at least three in \(T\). If \(n=2, \; n_3=0\) and \(T\) has two leaves.
Now suppose that every tree on fewer than \(n\) vertices has at least \((n_3+2)\) leaves, and consider a tree \(T\) on \(n \geq 3\) vertices. Since \(T\) is a tree on at least three vertices, it has a leaf \(u\). Then \(T'=T-u\) is a tree on \((n-1)\) vertices. By the induction hypothesis, we have 
\begin{equation}l(T') \geq n_3(T') + 2\end{equation}
\newpage\noindent Let \(v\) be the unique neighbour of \(u\) in \(T\). \(T\) is connected and has at least three vertices, so \(v\) has at least two neighbours in \(T\). This results in one of three cases:
\begin{enumerate}
    \item Suppose that \(v\) has exactly two neighbours in \(T\). Then \(n_3(T)=n_3(T')\) and \(l(T)=l(T')\). Hence,
    \begin{equation}
        l(T) = l(T') \geq n_3(T') + 2 = n_3(T) + 2
    \end{equation}
    \item Suppose that \(v\) has exactly three neighbours in \(T\). Then \(n_3(T) = (n_3(T') + 1)\) and \(l(T) = (l(T') + 1)\). Hence,
    \begin{equation}
        l(T) = l(T') + 1 \geq n_3(T') + 2 + 1 = n_3(T) - 1 + 2 + 1 = n_3(T) + 2
    \end{equation}
    \item Suppose that \(v\) has at least four neighbours in \(T\). Then \(n_3(T) = n_3(T')\) and \(l(T) = (l(T') + 1)\). Hence,
    \begin{equation}
        l(T) = l(T') + 1 \geq n_3(T') + 2 + 1 = n_3(T) + 2 + 1 \geq n_3(T) + 2
    \end{equation}
\end{enumerate}
In each case, \(l(T)\) is at least \(n_3+2\).}

\separate{5}\theorem{11.1.6}{Every tree is a bipartite graph.}{Let \(T\) be a tree and \(u\) be any one of its vertices. Place \(v\) in the set \(V_1\). For every vertex \(u \neq v\), there is a unique path from \(v\) to \(u\) in \(T\), consider the length of this path. If the length is odd, put \(u\) in \(V_2\), otherwise, place it in \(V_1\). \(V_1\) and \(V_2\) are disjoint and together make up \(V(T)\) and every edge of \(T\) will have an endpoint in both \(V_1\) and \(V_2\).}

\subsection{Isomorphism}

\definition{11.2.1}{Two graphs \(G(V,E)\) and \(G'(V',E')\) are \underline{isomorphic} if there exists a bijective function \(f:V \rightarrow V'\) such that for every \(u,v \in V\), we have \(uv \in E\) iff \(f(u)f(v) \in E'\). Then we write \(G \cong G'\).}

\separate{5}\definition{11.2.2}{Two rooted trees \(T_1(V_1,E_1,r_1)\) and \(T_2(V_2,E_2,r_2)\) are called \underline{isomorphic rooted trees} if there exists an isomorphism bijection \(f:V_1\rightarrow V_2\) such that \(f(r_1)=r_2\).}

\separate{5}\definition{11.2.3}{In a rooted tree with root \(r\), the \underline{level}, \(L(i)\) is the set of vertices at distance \(i\) from the root \(r\).}

\algorithmBlock{Label Vertex}{\(T\), a tree, and \(v\), a vertex from the tree.}{A label for vertex \(v\).}
{
    \If {\(v\text{ is a leaf of }T\)}
        \State \(label(v) \gets 10\)
    \Else
        \For {\(\text{every child } w \text{ of } v\)}
            \State \(l(w) \gets LabelVertex(T,w)\)
        \EndFor
        \State Sort the labels of the children decreasingly.
        \State \(label(v) \gets 1\)
    \EndIf
    
    \separate{3}\Return \(label(v)\)
}

\newpage\algorithmBlock{Labelled Tree Isomorphism}{Two pairs of trees and roots, \((T_i, r_i) \; i \in \{1,2\}\).}{YES or NO.}
{
    \State \(label(r_1) \gets LabelVertex(T_1, r_1)\)
    \State \(label(r_2) \gets LabelVertex(T_2, r_2)\)

    \separate{3}\If {\(label(r_1) = label(r_2)\)}
        
    \Return YES
    \Else 
    
    \Return NO
    \EndIf 
}

\noindent\theorem{11.2.1}{If \((T_1,r_1)\) is isomorphic to \((T_2,r_2)\) then \hyperref[alg:Labelled Tree Isomorphism]{Algorithm 4} returns YES.}
{We use induction on the number \(k\) of levels. When \(k=1\), the two trees have only their root (thus must be isomorphic), which gets label 10. Thus the algorithm returns YES. If two rooted trees are isomorphic and both have \(k\) levels then the algorithm returns YES on input. Let \((T_1,r_1)\) and \((T_2,r_2)\) be two isomorphic rooted trees, each with \((k+1)\) levels. 

Let \(f\) be the bijection between \(T_1\) and \(T_2\). Then \(r_2=f(r_1)\) and, for every child \(a_1\) of \(r_1\) there exists a child \(a_2\) of \(r_2\) such that the subtrees rooted at \(a_i\) are isomorphic, and \(a_2 = f(a_1)\). Since every such pair of subtrees are isomorphic, they both have \(k\) levels. Thus, by the induction hypothesis, \hyperref[alg:Label Vertex]{Algorithm 3} returns the same label for their roots. Since all these labels are the same, \hyperref[alg:Labelled Tree Isomorphism]{Algorithm 4} returns YES.}

\section{Minimum Spanning Trees}

\subsection{Introduction to Minimum Spanning Trees}

\definition{12.1.1}{The \underline{minimum spanning tree problem} tries to find a tree that spans all the vertices in a graph for the lowest possible cost. Minimum spanning trees have \((|V|-1)\) edges, have no cycles, and might not be unique.}

\separate{5}\definition{12.1.2}{\underline{Kruskal's algorithm} attempts to solve the problem by:
\begin{enumerate}
    \item Sort the edges by weight.
    \item Let \(A = \emptyset\)
    \item Consider the edges in increasing order of weight. For each edge \(e\), add \(e\) to \(A\) unless this would create a cycle.
\end{enumerate}
It has a running time of \(O(E\,log\,V)\).}

\separate{5}\definition{12.1.3}{\underline{Prim's algorithm} attempts to solve the problem by:
\begin{enumerate}
    \item Let \(U = \{u\}\) where \(u\) is some vertex chosen arbitrarily.
    \item Let \(A = \emptyset\)
    \item Until \(U\) contains all vertices: find the least-weight edge \(e\) that joins a vertex \(v\) in \(U\) to a vertex \(w\) not in \(U\) and add \(e\) to \(A\) and \(w\) to \(U\).
\end{enumerate}
It has a running time of \(O(V\,log\,V + E)\).}

\separate{5}\definition{12.1.4}{Let \(A\) be a subset of an MST. If the set \(A \cup \{(u,v)\}\) is also a subset of an MST, then \((u,v)\) is a \underline{safe} edge for \(A\).}

\separate{5}\definition{12.1.5}{A \underline{cut} \((S, V-S)\) of \(G=(V,E)\) is a partition of \(V\).}

\separate{5}\definition{12.1.6}{An edge \((u,v) \in E\) \underline{crosses} a cut \((S,V-S)\) if \(u \in S\) and \(v \in V-S\), or vice-versa.}

\separate{5}\definition{12.1.7}{An edge \((u,v) \in E\) is a \underline{light} edge crossing a cut if its weight is smallest among all the edges crossing the cut.}

\separate{5}\definition{12.1.8}{A cut \underline{respects} a set \(A\) of edges if no edge of \(A\) crosses the cut.}

\algorithmBlock{Generic MST Algorithm}{A graph \(G\).}{An MST.}
{
    \State \(A \gets \emptyset\)
    \While {\(A \text{ does not form a spanning tree}\)}
        \State Find an edge \((u,v)\) that is safe for \(A\)
        \State \(A \gets A \cup \{(u,v)\}\)
    \EndWhile
    
    \separate{3}\Return \(A\)
}

\subsection{Implementing MST Algorithms}

\algorithmBlock{Implementation of Kruskal's Algorithm}{The set of vertices \(V\) and the set of edges \(E\).}{An MST.}
{
    \State \(A \gets \emptyset\)
    \For {\(\text{Each vertex } v \in V\)}
        \State \(C(v) \gets {v}\)
    \EndFor
    \State Sort \(E\)
    \While {\(E \neq \emptyset\)}
        \State \(e \gets (u, v) \in E\) with min cost
        \If {\(C(u) \neq C(v)\)}
            \State \(A \gets (A+e)\)
            \For {\(\text{Each vertex } w \in C(u), w \in C(v)\)}
                \State Update \(C(w)\) with \(C(u) \cup C(v)\)
            \EndFor
        \EndIf
    \EndWhile
    
    \separate{3}\Return \(A\)
}

\noindent The sorting initially takes \(O(E\,log\,V)\) time, the union-find operations also take \(O(E\,log\,V)\) time, so the total running time is \(O(E\,log\,V)\).

\algorithmBlock{Implementation of Prim's Algorithm}{The set of vertices \(V\) and the set of edges \(E\).}{An MST.}
{
    \State \(U = \{u\}\)
    \State \(A \gets \emptyset\)
    \For {\(\text{Each vertex }v \in (V-u)\)}
        \State \(B(v)\) is the least-weight edge from \(v\) to \(U\)
    \EndFor
    \While {\(U \neq V\)}
        \State Choose \(v\) with minimum cost \(B(v)\)
        \State \(A \gets (A + e)\)
        \State \(U \gets (U + v)\)
        \State Update \(B\)
    \EndWhile
    
    \separate{3}\Return \(A\)
}

\noindent To initialise the algorithm, all edges are considered. It then iterates through the while loop once per vertex extracting the minimum cost edge and performing updates in \(O(log\,V)\) time. This gives an overall running time of \(O(V\,log\,V + E)\).

\newpage\part{References}

\section{Module Professors}

The below is accurate as of the 2025/26 academic year and will not be updated in subsequent years.
\begin{itemize}
    \item \hyperref[part:1.fds]{\textit{Fundamental Data Structures}} and \hyperref[part:4.gta]{\textit{Graph Theory and Algorithms}} are taught by David Kutner.
    \begin{itemize}
        \item \textbf{Office Location}: Room 2011, MCS Building
        \item \textbf{Website}: https://www.durham.ac.uk/staff/david-c-kutner/
        \item \textbf{Preferred Contact Methods}: david.c.kutner@durham.ac.uk
        \item \textbf{Pronouns}: he, him
    \end{itemize}
    
    \item \hyperref[part:2.ans]{\textit{Asymptotic Notation and Sorting}} is taught by Thomas Erlebach.
    \begin{itemize}
        \item \textbf{Office Location}: Room 2004, MCS Building
        \item \textbf{Website}: https://www.durham.ac.uk/staff/thomas-erlebach/
        \item \textbf{Preferred Contact Methods}: thomas.erlebach@durham.ac.uk
        \item \textbf{Pronouns}: he, him
    \end{itemize}
    
    \item \hyperref[part:3.sds]{\textit{Selection and Data Structures}} is taught by Anish Jindal.
    \begin{itemize}
        \item \textbf{Office Location}: Room 2065, MCS Building
        \item \textbf{Website}: https://www.durham.ac.uk/staff/anish-jindal/
        \item \textbf{Preferred Contact Methods}: anish.jindal@durham.ac.uk
        \item \textbf{Pronouns}: he, him
    \end{itemize}
\end{itemize}

\section{Cited Works}

Erickson, J. (2019). Algorithms. 1st ed. [online] Self-Published, p.71. \\ Available at: http://jeffe.cs.illinois.edu/teaching/algorithms/ [Accessed 25 Mar. 2026].\label{cit:1.algorithms}

\separate{5}Rosen, K.H. (2003). Instructor's Resource Guide for Discrete Mathematics and Its Applications. \\ Boston: Mcgraw-Hill, pp.634-636.\label{cit:2.eulerianCircuits}

\end{document}

%%%